\section{\label{sec:rgeqns}重正化群方程}

我们考虑标量算符与 $N$ 个外部未涂抹标量场耦合的矩阵元。通过考察恰当选取的格林函数的标度依赖，
可以导出 OPE 威尔逊系数的重正化群方程 \cite{Collins:1984rdg}。$N+2$ 个外部标量场的格林函数
满足重正化群方程
\begin{equation}\label{eq:rg1}
\left[\mu\frac{\mathrm{d}}{\mathrm{d} \mu}+\left(N+2\right)\gamma
\right]\big\langle \Omega |
\widetilde{\phi}(p_1) \ldots\widetilde{\phi}(p_{N+2}) | \Omega \big\rangle = 0,
\end{equation}
而重正化算符 $\phi_{\mathrm{R}}^2(0)$ 与 $N$ 个外部标量场耦合的格林函数则满足
\begin{equation}\label{eq:rg2}
\left[\mu\frac{\mathrm{d}}{\mathrm{d} \mu}-\gamma_m + N\gamma
\right]\big\langle \Omega |  \phi_{\mathrm{R}}^2(0)
\widetilde{\phi}(p_1) \ldots\widetilde{\phi}(p_N) | \Omega \big\rangle = 0.
\end{equation}
这里，标量场论的重正化群算符为
\begin{equation}
\mu\frac{\mathrm{d}}{\mathrm{d} \mu} = \mu\frac{\partial}{\partial
\mu}\bigg|_{\lambda,m_{\mathrm{R}}} + \beta
\frac{\partial}{\partial \lambda}\bigg|_{\mu,m_{\mathrm{R}}} - \gamma_m
m_{\mathrm{R}}\frac{\partial}{\partial m_{\mathrm{R}}}\bigg|_{\mu,\lambda}
\end{equation}
各系数为 \cite{Kleinert:2001ax}
\begin{align}
\beta^{\overline{MS}}(\lambda) = {} & \mu
\frac{\mathrm{d}\lambda}{\mathrm{d}\mu}
=\frac{3\lambda^2}{(16\pi^2)^2}+{\cal O}(\lambda^3), \\
\gamma_m^{\overline{MS}}(\lambda) = {} &
-\frac{\mu}{2}\frac{\mathrm{d}\log(m_{\mathrm{R}})}{\mathrm{d}\mu}
\nonumber \\
= {} &  -\frac{\lambda}{2(16\pi^2)}+
\frac{5\lambda^2}{12(16\pi^2)^2}+{\cal O}(\lambda^3), \label{eq:gammam}\\
\gamma^{\overline{MS}}(\lambda) = {} & \frac{\mu}{2}
\frac{\mathrm{d}\log(Z_\phi)}{\mathrm{d}\mu}
=  \frac{\lambda^2}{12(16\pi^2)^2}+{\cal O}(\lambda^3).
\end{align}
一般情形下这些系数依赖于质量、重正化尺度以及重正化的耦合常数；但在与质量无关的重正化方案
（如 $\overline{MS}$ 方案）中，
它们仅通过重正化耦合常数依赖于重正化尺度。对 $O(N)$ 对称理论（由 $N$
重态 $\Phi = \{\phi_1,\ldots,\phi_N\}$ 给出）而言，
这些函数已知到五圈 \cite{Derkachov:1997pf,Kleinert:2001ax}。

再次回到两点函数 OPE 的例子、即式~\eqref{eq:opeint}：
把算符与两个外部标量场耦合并利用式~\eqref{eq:rg1} 和 \eqref{eq:rg2}
（更多细节可参见例如 Ref.~\cite{Collins:1984rdg}），可以导出威尔逊系数
$c_{\phi^2}$ 的重正化群方程：
\begin{equation}
\left[\mu\frac{\mathrm{d}}{\mathrm{d} \mu}
 + 2\big(\gamma +\gamma_m\big)\right]c_{\phi^2} = 0.
\end{equation}

正如我们所预期的，威尔逊系数 $c_{\phi^2}$ 的反常维数等于
$\phi(x)\phi(0)$ 的反常维数与 $[\phi^2(0)]_{\mathrm{R}}$ 反常维数之差。

\subsection{sOPE 威尔逊系数的重正化群方程}

流时间 $\tau$ 为问题引入了一个新尺度。原则上可以把流时间仅仅视为另一个外部尺度，
与重正化尺度 $\mu$ 无关。此时，自然要修改重正化群方程，以顾及 $\mu$ 与
$\tau$ 同时变化时格林函数的改变：
\begin{equation}
\mu\frac{\mathrm{d}}{\mathrm{d} \mu}\rightarrow
\mu\frac{\mathrm{d}}{\mathrm{d}
\mu}- 2\tau \frac{\mathrm{d}}{\mathrm{d} \tau}.
\end{equation}
我们对微分算符的选择受每个尺度（$\mu$ 与 $\tau$）的质量量纲约束，但并非唯一。特别地，也可以选择算符
$\mu\mathrm{d}/\mathrm{d}
\mu + \tau \mathrm{d}/\mathrm{d} \tau$。这种自由既不影响我们论证的逻辑，也不影响我们的结论，
因为不同的约定可以被吸收进——例如——式~(58)
中重正化参数与反常维数的定义之中。

此处值得对问题中的两个尺度 $\mu$ 与 $\tau$ 加以评论。在 DIS 中，
相应 OPE 的时间空间隔 $x$ 与重正化尺度 $\mu$ 原则上是两个不同的尺度。时间空间隔由某个（或某组）
DIS 实验的动量转移之倒数给出；而重正化尺度是理论上的选择，最终物理量不应依赖于它。
选取 $\mu = 1/x$ 通常很方便，但并非严格必需。

流时间与重正化尺度之间的关系与之类似，这两个尺度也是彼此不同的。对 sOPE 而言，
流时间可以被视为仅是由某种特定格点``实验''施加的外部尺度，而重正化尺度则是便利的理论选择。
我们将看到，把这两个尺度联系起来有助于把双尺度问题约化为单尺度问题，但这在形式上并非必要。
因此，在下文的分析中，流时间与重正化尺度应被理解为完全独立的尺度。

再次考虑两点函数的 sOPE（式~\eqref{eq:sope2pt}），我们按照与上文类似的手续确定涂抹威尔逊系数
$d_{\rho^2}$ 的重正化群方程。

我们假定处于小流时间区域；这使我们能把零流时间与非零流时间的算符联系起来
\cite{Makino:2014taa,Suzuki:2013gza,Luscher:2011bx}：
\begin{equation}\label{eq:smallt}
[\phi^2(0)]_{\mathrm{R}}  =   {\cal
Z}_{\rho^2}(\tau,\mu)\rho^2(\tau,0) + {\cal O}(\tau).
\end{equation}
该系数满足
\begin{equation}\label{eq:zgamma}
\mu\frac{\mathrm{d}}{\mathrm{d}\mu}\log({\cal
Z}_{\rho^2}(\tau,\mu^2)) = 2\gamma_m,
\end{equation}
且其单圈形式为
\begin{equation}
{\cal Z}_{\rho^2}(\tau,\mu^2)=
1-\frac{\lambda}{32\pi^2}\left[1+\gamma_{\;\mathrm{E}}+\log(2\tau
\mu^2)\right]+{\cal
O}(\lambda^2,\tau).
\end{equation}

我们把重正化群算符
\begin{equation}
\mu\frac{\mathrm{d}}{\mathrm{d} \mu}
-2\tau\frac{\mathrm{d}}{\mathrm{d} \tau} + \left(N+2\right)\gamma
\end{equation}
作用于式~\eqref{eq:smallt} 中诸算符与 $N$ 个外部标量场耦合的矩阵元，得到
\begin{equation}\label{eq:srge}
\left[\mu\frac{\mathrm{d}}{\mathrm{d}
\mu} -2\tau\frac{\mathrm{d}}{\mathrm{d}
\tau} +2(\zeta_{\rho^2}-\gamma)\right]d_{\rho^2} = {\cal O}(\tau).
\end{equation}
其中
\begin{equation}
\zeta_{\rho^2} = \tau\frac{\mathrm{d}}{\mathrm{d}\tau}\log({\cal
Z}_{\rho^2}(\tau,\mu^2))\label{eq:zetadef}
\end{equation}
是与算符 $\rho^2(\tau,0)$ 流时间依赖相伴的一个反常维数。$\rho^2(\tau,0)$
与 $N$ 个外场耦合的相应矩阵元的重正化群方程为
\begin{align}\label{eq:rhotau}
\bigg[\mu\frac{\mathrm{d}}{\mathrm{d}
\mu}{} & -2\tau\frac{\mathrm{d}}{\mathrm{d}
\tau} +2\zeta_{\rho^2}+N\gamma\bigg] \nonumber\\
{} & \times\big\langle \Omega
|\rho^2(\tau,0)\widetilde{\phi}(p_1)\ldots\widetilde{\phi}(p_N)
| \Omega \big\rangle = {\cal O}(\tau).
\end{align}
我们注意到，该方程仅在流时间小于外部粒子动量（对 DIS 而言为强子尺度的量级）时才成立。

如果现在要求涂抹尺度 $\tau$ 与重正化尺度 $\mu$ 的倒数互成比例，即
$\tau = b/\mu^2$（$b$ 为实数），那么该重正化群方程变为
\begin{align}\label{eq:rhotau_mutau}
{} &
\bigg[2\mu\frac{\mathrm{d}}{\mathrm{d}
\mu} +2\zeta_{\rho^2}+N\gamma\bigg] \nonumber\\
{} &\qquad \times\big\langle \Omega
|\rho^2\left( b/\mu^2,0\right)\widetilde{\phi}(p_1)\ldots\widetilde{\phi
} (p_N)
| \Omega \big\rangle = {\cal O}(b).
\end{align}
这个重正化群方程是非微扰步标度方法
\cite{Monahan:2013lwa,Luscher:1991oxn} 的出发点：该方法把非微扰矩阵元演化到高标度，
在那里可以将其与微扰的涂抹威尔逊系数相结合。这里，重正化群方程在小流时间极限下、
换言之在 $b\ll 1$ 时成立。这一约束自动保证流时间也小于任何强子长度尺度、
$\Lambda_{\mathrm{QCD}} \ll \mu^2 \ll 1/\tau$，并且对实用的步标度方法一般成立
\cite{Monahan:2013lwa,Luscher:1991oxn}。我们目前正在 QCD 中研究这一途径。
